% !TEX program = latexmk
% !TEX options = -pdf -synctex=1 -interaction=nonstopmode -file-line-error "%DOC%"
%
% ==============================================================================
% - Maximum 20 pages (excluding Title Page, Abstract, and References)
% - 1.5 line spacing
% - 1-inch margins on all sides
% - 11pt minimum font size (Times New Roman / MathPTMX utilized here)
% - Non-bold section headers (standard scholarly/STS submission style)
% ==============================================================================

\documentclass[11pt, letterpaper]{article}

% --- Core Formatting & Page Layout ---
\usepackage[margin=1in]{geometry}
\usepackage{setspace}
\onehalfspacing

% --- Typography & Math Setup ---
\usepackage{authblk}
% \usepackage{newtxtext,newtxmath}                  % Times font family
\usepackage{mathtools}
\usepackage{amsmath, amsfonts, amssymb} % Advanced mathematical modeling & matrix operations
\usepackage{graphicx}                  % Figure inclusion
\usepackage{booktabs}                  % Professional table formatting
\usepackage{cite}                      % Standard numerical citation management
\usepackage{hyperref}                  % Hyperlinks and metadata
\usepackage{dsfont}

% --- Non-Bold Section Title Customization ---
\usepackage{titlesec}

% Section: Large, Normal Weight (Non-Bold)
% \titleformat{\section}
%   {\normalfont\Large\fontfamily{ptm}\selectfont}{\thesection}{1em}{}

% % Subsection: Medium, Normal Weight (Non-Bold)
% \titleformat{\subsection}
%   {\normalfont\large\fontfamily{ptm}\selectfont}{\thesubsection}{1em}{}

% % Subsubsection: Normal Size, Normal Weight (Non-Bold)
% \titleformat{\subsubsection}
%   {\normalfont\normalsize\fontfamily{ptm}\selectfont}{\thesubsubsection}{1em}{}

% % Adjust spacing around section titles
% \titlespacing*{\section}{0pt}{12pt}{6pt}
% \titlespacing*{\subsection}{0pt}{10pt}{4pt}
% \titlespacing*{\subsubsection}{0pt}{8pt}{2pt}

% --- Hyperlink Setup ---
\hypersetup{
    colorlinks=true,
    linkcolor=black,
    filecolor=black,
    urlcolor=blue,
    citecolor=black
}

% --- Code & Algorithmic Listing Setup ---
\usepackage{listings}
\usepackage{xcolor}

\definecolor{codegreen}{rgb}{0,0.6,0}
\definecolor{codegray}{rgb}{0.5,0.5,0.5}
\definecolor{codepurple}{rgb}{0.58,0,0.82}
\definecolor{backcolour}{rgb}{0.98,0.98,0.98}

% \lstdefinestyle{stsstyle}{
%     backgroundcolor=\color{backcolour},
%     commentstyle=\color{codegreen},
%     keywordstyle=\color{blue},
%     numberstyle=\tiny\color{codegray},
%     stringstyle=\color{codepurple},
%     basicstyle=\ttfamily\footnotesize,
%     breakatwhitespace=false,
%     breaklines=true,
%     captionpos=b,
%     keepspaces=true,
%     numbers=left,
%     numbersep=5pt,
%     showspaces=false,
%     showstringspaces=false,
%     showtabs=false,
%     tabsize=4
% }
% \lstset{style=stsstyle}

\title{Penalized Portfolio Optimization}
\author{Victor Ding}
\affil{New Canaan High School}
\date{\today}
\begin{document}

\maketitle

\newpage
\setcounter{page}{1} % Page 1 starts with the Abstract
\pagestyle{plain}

\begin{abstract}
\noindent
This research reviews practical considerations of long-short portfolio optimization problems. It studies a few penalty forms to
address the limitation of pure constraints.
\end{abstract}
\newpage
\section{Background of Portfolio Construction}

\subsection{Principles of Portfolio Construction}

Modern portfolio construction involves optimizing an objective function defined over a vector of $n$ security positions
$w \in \mathds{R}^n$, where positive (negative) values represent long (short) positions. An example of a classic objective
function is Markowitz's Mean-Variance framework:
\begin{equation}
    \label{mean-variance_objective}
    \underset{w}{\min} \left[\frac{\lambda}{2} w' \Sigma w - r' w\right],
\end{equation}
where $r$ is a vector of expected returns, $\lambda$ is an exogenous risk-aversion parameter, and $\Sigma$ is a covariance
matrix of asset returns. The design is driven by the desire to strike a balance between risks and payoffs.

Objective functions can take other forms, such as $L_2$ norm of the deviation from a target vector $\tilde{w} \in
\mathds{R}^n$. A given vector $\tilde{w}$ can be a benchmark portfolio to track, or a target portfolio driven by some
forecasting models. It is often desirable to mimic such a portfolio as closely as possible. One reasonable approach is to
minimize this objective:
\begin{equation*}
	\underset{w}{\min} \lVert w - \tilde{w}\rVert_2^2.
\end{equation*}
A non-negative diagonal weighting scheme matrix $\Gamma$ can also be incorporated to generalize the above objective function further:
\begin{equation}
  \label{gamma_matrix}
	\underset{w}{\min}\left[\frac{1}{2} (w - \tilde{w})'\Gamma(w - \tilde{w})\right] \equiv
  \underset{w}{\min}\left[\frac{1}{2} w'\Gamma w - \tilde{w}'\Gamma w \right].
\end{equation}
Mathematically, this problem is a special case of \eqref{mean-variance_objective}. It will be elaborated upon in section
\ref{target_objective_section}.

A majority of portfolio construction problems can be formulated in certain quadratic forms like the above. Typical constraints
include position limits, risk exposure and turnovers, etc.

Objective \eqref{mean-variance_objective} is convex if $\Sigma$ is positive semi-definite and if the constraints
maintain a convex feasible set. Objective \eqref{gamma_matrix} acts similarly, placing both objectives within the reach of
Disciplined Convex Programming (DCP). DCP problems can be easily solved with DCP solvers.

% \subsection{Research gap related to constraints and penalties}
% COMPLETE LATER
% Penalties occupy a middle-ground between an objective function and constraints; instead of placing a
% hard boundary on portfolios, they discourage going past a certain threshold continuously. If an objective function
% is minimized, penalties come in the form of added terms.


\section{Implementation Framework}

Let us focus on minimizing the family of objectives with a quadratic term $q(w)$, which captures the conventional Mean-Variance
problem \eqref{mean-variance_objective} and target tracking problem \eqref{gamma_matrix}.

MOVE LATER
Constraints model numerous types of boundaries, including risk modeling. This includes factor exposure bounds,
sector or industry neutrality, and tracking error limits.

Inputs such as return estimates, risk model coefficients, and target forecasts in $\tilde{w}$ are all estimates subject to
statistical errors. Thus, a problem subject to overly permissive constraints is susceptible to concentrated exposure on assets
with potentially poor signals. Constraints like turnover and position limits arise from practical considerations and also serve as a
form of regularization on model solutions.

On the other hand, we look to avoid overly-tight constraints. While constraints have their purposes, each one shrinks the feasible set of
potential solutions. Certain constraints may result in solutions that reflect little of the objective. Therefore, constraints should be set
such that they do not "drive" the solution; the solution should be dictated in part by the objective function.

While potential payoffs are often uncertain, transaction costs are incurred upon execution. Turnover budgets and individual
turnovers are often imposed, so a solution is discouraged from chasing marginal benefits with excessive transaction costs.

REFORMAT
Linear equalities, inequalities, and second-order cones define half-spaces, hyperplanes, and second-order cones respectively.
These all define convex sets. Optimization works better
when the feasible set is convex. Because of this, we show a preference for preserving convexity of the overall
feasible region. For constraints that are non-smooth but inherently convex (such as absolute values or piecewise functions), we
exercise exact epigraph reformulations with auxiliary variables. For constraints that are truly non-convex, we employ convex
relaxations to approximate the feasible space while maintaining computational tractability.

\subsection{Constraints}
Position limits can be represented as
\begin{align*}
    \underline{w} & \leq w \leq \overline{w},
\end{align*}
where constant vectors $\underline{w} \in \mathds{R}^n$ and $\overline{w} \in \mathds{R}^n$ are upper and lower bounds, respectively.

Similarly, limiting each individual turnover is represented by
\begin{equation*}
    |w - w^*| \leq t,
\end{equation*}
where constant $t \in \mathds{R}^n$ defines individual turnover bounds. Turnover bounds are essentially a different set of
position limits because $w^*$ is a constant vector. They can be merged together.

We define an attribute as a characteristic of an asset, which encompasses risk modeling, equity style factors, and more. Risk
modeling is often adopted in portfolio construction to limit risk exposure, commonly in sectors and industries. Equity style
factors (growth, momentum, value, etc.) are also attributes, and can be used in risk modeling. These detail different
predictions on how certain assets can outperform others.

Assets can also have exposure to multiple attributes. For example, although a small tech startup will have exposure in just a technology sector, a company like
Alphabet operates in several. Companies can also have negative exposure in attribtues if their growths have an inverse relationship with attributes.

We express $n$ assets' exposure in $m$ attributes with an $m\times n$ matrix $\mathbf{S}$. For example, a potential $\mathbf{S}$
for $n$ assets across various attributes is shown in Table \ref{example_attribute_matrix}.  The total portfolio exposures to the
$m$ factors are $\mathbf{S}w$. The exposure limits can be represented as
\begin{equation}
    \underline{s} \leq \mathbf{S}w \leq \overline{s}
\end{equation}

\begin{table}[t]
\label{example_attribute_matrix}
\centering
    \caption{Example Factor Loading Matrix.}
    \vspace{0.3cm}
    \begin{tabular}{cr *{5}{c}}
        \toprule
        & & \multicolumn{5}{c}{\textbf{Assets}} \\
        \cmidrule(lr){3-7}
        % Replace the 3rd column with dots, and 4th column with the n-value
        & & \textbf{1} & \textbf{2} & \textbf{3} & \dots & \textbf{$n$} \\
        \cmidrule(lr){1-7}
        & Finance & 0.10 & 0.40 & 0 &\dots & \dots \\
        & Energy & 0.13 & 0.20 & \dots & \dots & \dots \\
        & Communication & 0.06 & 0.04 & \dots & \dots & \dots \\
        \cmidrule(lr){1-7}
        & \dots & \dots & \dots & \dots & \dots & \dots \\
        & Growth & 0.70 & 0.30 & & \dots & \dots \\
        & Momentum & -1.00 & -0.30 & & \dots & \dots \\
        & \dots & \dots & \dots & \dots & \dots & \dots \\
        \bottomrule
    \end{tabular}
\end{table}

Net exposure is a special example of an attribute that can be included as a row of all ones in $\mathbf{S}$. It is equal to the
value of long exposure positions minus the value of short exposure positions. The net exposure of a portfolio is $e' w$, where
$e$ is a vector of ones.

A constant upper bound $T > 0$ may be imposed on total turnover as
\begin{equation}
  \label{unedited_total_turnover}
	\lVert w^* - w\rVert_1 \leq T.
\end{equation}
This introduces a total budget over the sum of individual turnovers. In practice, individual turnover limits in $t$ are often
set to such that $T \leq \sum_i t_i$; otherwise, the total turnover cap will become redundant. It can be further generalized
with a weighting vector $c > 0$ such that
\begin{equation*}
	c' |w^* - w| \leq T.
\end{equation*}
The vector $c$ and $T$ may represent some transaction cost estimates and a budge limit, respectively.

The constraint \eqref{unedited_total_turnover} is not directly supported by many solvers. Although convex, it is not affine, and
is not differentiable at $w_i = w_i^*$. To alleviate this, slack variables can be created to reformulate the constraint. This
operation is explained in detail in section \ref{total_turnover_section}.

Gross exposure, $\lvert w\rVert_1$, is another measure often involved in some constraints, as firms often wish to keep a
specific amount of capital invested in a portfolio. It is common to limit the gross exposure (leverage) by an upper bound $G$ as
\begin{equation*}
	\lvert w\rVert_1 \leq G.
\end{equation*}
This is still a convex constraint like that on total turnovers. But the optimal solution may find it cheaper to settle at a
gross exposure far less than $G$, effectively underinvesting with available capital. It is thus attractive to place a lower
bound $\underline{G}$ on the gross exposure as well. When $G=\underline{G}$, the constraint becomes an equality one.
\begin{equation}
  \label{explicit_gross_exposure}
  \underline{G} \leq	\lvert w\rVert_1 \leq G.
\end{equation}
In long-only or short-only portfolios, absolute value signs in \eqref{explicit_gross_exposure} can be trivially removed so that
it becomes a convex constraint. But the presence of a lower bound or equality constraints in long-short portfolios makes the
feasible set non-convexity. This poses a problem for most solvers in its current form. We will need to address this
non-convexity in section \ref{gross_exposure_section}.

\subsection{Karush-Kuhn-Tucker and Global Optimality Conditions}
\label{KKT_section}
Because the portfolio construction problems considered henceforth are convex, their solutions are
exact and can be calculated via DCP. Suppose a quadratic objective is bounded by affine constraints that
create a convex feasible set. The Karush-Kuhn-Tucker (KKT) conditions \cite{boyd2004KKT} define the framework for determining
optimality. For a convex objective function with convex inequality and affine equality constraints, a point
that satisfies the conditions is guaranteed to be a global minimum. Using a convex quadratic form to satisfy these conditions
calls for two conditions: a positive semi-definite Hessian matrix, which indicates convexity, and a
constraint qualification, which we will fulfill through using affine constraints.

In its current state, the objective function satisfies these conditions as it contains purely quadratic terms, and almost all constraints
implemented thus far satisfy, the exception being gross exposure.
\subsection{Treatment of Total Turnover}
\label{total_turnover_section}
Many solvers do not natively accept constraints with $L1$ norms like \eqref{unedited_total_turnover}. A slack vector
$y \in \mathds{R}^n$ is introduced to replace it with
\begin{equation}\label{eq:tvr-slack}
	|w^* - w| \leq y, \quad e'y \leq T.
\end{equation}
Since increasing $y_i$ will only make this condition more difficult to satisfy and doesn't affect the objective value, $y_i$
will be forced downwards toward the tightest feasible value $y_i = |w_i^* - w_i|$. In doing so, \eqref{unedited_total_turnover}
is replaced with strictly linear expressions. At the optimal value, it must satisfy
\begin{equation*}
    e' y = \lVert w^* - w\rVert_1 \leq T.
\end{equation*}
These reformulated constraints are strictly linear.


\subsection{Difficulties of Gross Exposure}
\label{gross_exposure_section}
A set $\mathcal{C} \subset \mathds{R}^n$ is convex if and only if for any two points $a$ and $b$ in $\mathcal{C}$ and every
$\theta \in [0,1]$ the following is also true:
\begin{equation}
    \label{convexity_definition}
	\theta a + b (1 - \theta) \in \mathcal{C}.
\end{equation}

We will prove that for fixed $G > 0$, the set $\mathcal{S} = \{w\in \mathds{R}^n:\lVert w\rVert_1 = G\}$ is not convex.
This argument does not restrict itself to properties of the $L_1$ norm. Rather, it uses two axioms every norm satisfies:
\begin{enumerate}
	\item Absolute Homogeneity: $\lVert\alpha w\rVert = |\alpha| \lVert w\rVert$ for any scalar $\alpha$.
	\item Positive definiteness: $\lVert w\rVert = 0 \Leftrightarrow w = 0$.
\end{enumerate}
Take any $w \in \mathcal{S}$. By homogeneity with $\alpha = -1$,
\begin{equation*}
	\lVert-w\rVert_1 = | -1|\cdot \lVert w\rVert_1 = \lVert w\rVert_1 = G,
\end{equation*}
so $-w \in \mathcal{S}$ as well. Now take $\theta = \frac{1}{2}$ and form \eqref{convexity_definition} of $w$ and $-w$:
\begin{equation*}
	\theta w + (1-\theta)(-w) = 0.
\end{equation*}
By positive definiteness, $\lVert 0\rVert_1 = 0$. Since $G>0$ by assumption, $0 \neq G$, so
\begin{equation*}
	\theta w + (1-\theta)(-w) = 0 \notin \mathcal{S}.
\end{equation*}
We have shown that $\theta w + (1 - \theta) (-w) \notin \mathcal{S}$ for both $w$ and $-w$ in $\mathcal{S}$.
By \eqref{convexity_definition}, the set $\mathcal{S}$ is not convex.

The non-convexity arises from the uncertain signs of the optimal solution, which are not known prior to the solution
process. Even Mixed-Integer Convex Programming (MICP) will have difficulty finding a global optimal solution in a timely and
resource-efficient manner when $n$ is particularly large.

One way to overcome it is forcing the sign of $\{w_i\}$ according to some rules prior to setting up the gross exposure
constraint. Let $G^+ > 0$ be the desired long exposure, and $G^- < 0$ be the negative counterpart. The total gross exposure
becomes $G^+ - G^- = G$. We propose to match the sign of $w$ to those of $\tilde{w}$ for all $i$ such that
\begin{alignat*}{2}
	\sum_{i\in\mathcal{W}^+} {w_i} &= G^+,\quad & \mathcal{W}^+=\{i|\tilde{w}_i > 0\},\\
	\sum_{i\in\mathcal{W}^-} {w_i} &= G^-,\quad & \mathcal{W}^-=\{i|\tilde{w}_i < 0\},\\
	\sum_{i\in\mathcal{W}^0} {w_i} &= 0, & \mathcal{W}^0=\{i|\tilde{w}_i = 0\}.
\end{alignat*}
The price is that the feasible set is effectively shrunken, and consequently will not guarantee to yield the global optimal
solution of the original problem.

There are more ways to rid of non-convexity. Other local search methods may improve upon this approach:
\begin{itemize}
\item Solve the original problem with inequality, and use the signs of the solution vector to derive $\mathcal{W}^+$,
	$\mathcal{W}^-$, and $\mathcal{W}^0$.
\item After one iteration, flip the signs of whose securities $|w_i|$ is within a small band $\epsilon$ to update
  $\mathcal{W}^+$, $\mathcal{W}^-$, and $\mathcal{W}^0$. Solve again. If there is improvement, increase $\epsilon$ and repeat.
  If not, stop.
\item Solve the original problem with inequality, and scale all the weights such that $\lVert w\rVert = G$.  Feed the results
  back to solve again.
\end{itemize}

\section{Hard and Soft Caps of Transaction Costs}
\label{additional_gating}

Hard constraints that bound turnover and other portfolio characteristics have certain downsides. Problems should be set such that they are
reflective of practical and risk limits, rather than impose constraints that dictate optimality.
Violating a hard constraint by a negligible amount and by a comparatively huge amount have the exact same effect, meaning rounding
errors, tiny violations, and genuine infeasibility are treated the same. Furthermore, if for some $i$, $|w_i^* - \tilde{w}_i|$ is significantly
greater than $t_i$, there is a considerable chance that the stock should be considerably skewed in the direction of $\tilde{w}_i$. In those cases,
it is desirable to permit certain violations of hard constraints. To do this, we introduce soft constraints.

Soft constraints allow such violations while discouraging large ones by adding a penalty to the objective. We will consider three
distinct types of soft constraints: an alpha penalty, an Elastic Net, and a Berhu penalty.

\subsection{Soft Cap of Total Turnover}
The trading from $w^*$ to $w$ incurs transaction costs. Naturally, this requires balancing objective and turnover to avoid
paying deterministic upfront costs on insignificant or uncertain transactions. Imposing a hard budge limit on total turnover,
however, handicaps the portfolio construction on achieving its objective.

The change from $w^*$ to $w$ incurs transaction costs. Naturally, this requires balancing tracking error and turnover to avoid
paying fixed costs on insignificant transactions. Using a concrete turnover bound that doesn't discourage/encourage different
magnitudes of turnover thus has downsides. We wish to induce a degree of sparsity in a portfolio's turnovers, meaning making
fewer, concentrated trades rather than making several tiny trades that will not have sizable effect on portfolio output, but
incur fixed costs. In other words, we should increase the number of unchanged assets between optimization. A more aggressive
approach to execute this involves setting a certain threshold $\left |w^*_i - \tilde{w}_i \right|$ past which an asset will be
forced to remain unchanged. However, this hard constraint leaves no room for exceptions.

In addition to a hard upper bound $T$ on total turnover $\lVert w - w^*\rVert_1$ or as its replacement, we can introduce a slack
variable $z > 0$ to relax the hard bound such that
\begin{equation*}
	\lVert w - w^*\rVert_1 \leq \tilde{T} + z,
\end{equation*}
and add a penalty term to the objective function $q(w)$,
\begin{equation*}
	p(z) = \alpha z^2.
\end{equation*}
The exogenous parameter $\alpha > 0$ determines the cost of going above $\tilde{T}$. The variable $z$ will be optimized along
with $w$. The quadratic term has the benefit of having a comparable magnitude with many objective function forms of $q(w)$ such
that the choice of $\alpha$ can be more intuitive. It is completely reasonable to replace it with $\alpha |z|$. The presence of
an absolute value poses no difficult here given that $z \geq 0$.

This formulation is very intuitive. The total turnover is allowed to exceed $\tilde{T}$, but pays a penalty to do so. If
turnover is less than $\tilde{T}$, no penalty is added. The larger the value of $\alpha$, the more the final turnover is
discouraged from exceeding $T$, and vice versa. This ``soft cap'' is not a rigid hard bound, allowing some room to trade more
when situation dictates that. When both a ``hard cap'' and a ``soft cap'' are present, we generally have $\tilde{T} < T$ such
that $T$ serves as a final safeguard on the total turnover.

\section{Regularization of Noisy Trades}

While it is a common practice to cap the total turnover $\lVert w - w^*\rVert_1$ or limit some weighted transaction costs
$e'\Gamma |w - w^*|$ by a budget, it is also desirable to reduce the amount of inconsequential small trades. They are either
results of some computational or statistical noises, or are simply too insignificant to justify the execution or handling
costs.

We wish to induce a degree of sparsity in a portfolio's turnovers, eliminating noisy trades. A common approach to achieve it
involves setting a certain threshold for $\left |w^*_i - w_i \right|$ below which those $\{w_i\}$ will be forced to stay at
$\{w_i^*\}$, meaning no trade allowed. However, this hard rule may be too rigid, and the threshold setting also seems rather
arbitrary.

This problem is very similar to variable selection in various regression frameworks where coefficients of some variables are
shrunk towards zero to yield more parsimonious models. It is thus not surprising that similar techniques can be employed to
regularize $w - w^*$.
\subsection{Elastic Net}

Ridge Regression \cite{hoerl1970ridge} softens rigid parameters through adding an $L_2$ norm scaled by a multiplier. This allows a soft constraint while
simultaneously discouraging it past a certain point. Lasso regression \cite{tibshirani1996lasso} modifies the $L_2$ norm to instead be an $L_1$ norm.
The Elastic Net \cite{zou2005elasticnet} uses both Lasso and Ridge simultaneously. This concept can be applied to portfolio optimization,
penalizing noise trades:
\begin{equation}
	\label{elastic_net_objective}
	\gamma\left[\beta \lVert w-w^*\rVert_1 + \frac{1 - \beta}{2} \lVert w-w^*\rVert_2^2\right],
\end{equation}
where $\beta \geq 0$ and $\gamma > 0$ are adjustable exogenous parameters. The $L_1$ penalty promotes sparsity. It is a
mechanism that drives small weight turnovers to $0$, reducing the number of assets to trade. The $L_1$ multiplier $\beta$
governs the degree of sparsity on the portfolio. A larger $\beta$ reflects a higher marginal cost of initiating a new
trade. This shrinks the amount of sensible trades, forcing the optimizer to only execute trades where expected returns outweigh
a fixed transaction cost. The parameter $\gamma$ controls the overall scale of this penalty term relative to the objective
function $q(w)$.

The ${L_2}$ penalty promotes smoothness. It prevents a single element from changing drastically. Hence, turnover among
correlated assets is distributed more evenly than with the absolute value penalty. For example, the $L_2$ penalty resulting from
two turnovers $0.2$ and $0.2$ is smaller than the penalty resulting from $0.1$ and $0.3$, even though the total turnover is the
same. Thus, the $L_2$ turnover results in spreading turnover more evenly across assets. The $\beta$ parameters governs relative
strength of penalties between the two terms.

A penalty term \eqref{elastic_net_objective} and an alpha gating quadratic term can be added to a base quadratic objective
$q(w)$.
Each term in the objective is either an $L_1$ norm, $L_2$ norm, or a quadratic function. If all scalars $\alpha$, $\beta$,
and $\gamma$ are positive (which they should be, otherwise penalties will not work as intended), the entire objective
function is strictly convex. This ensures that any optimal solution is a global optimum. Note that an $L_1$ norm in
the objective function is permissible, while an $L_1$ norm equality constraint is not convex.

The greater $\gamma$ is the more trades will be pushed towards zero. Since both $\beta$ and $\gamma$ are non-negative, this
penalty term will preserve strictly convexity in a convex objective function. This guarantees that any local optimal solution is
also the global optimum.

% ADD OTHER FIELDS THAT UTILIZE ELASTIC FIELDS
\subsection{Berhu's (Reverse Huber's) Loss}
Elastic Net penalties subject portfolio weights to both $L_1$ and $L_2$ penalties, regardless of the magnitude of an asset's turnover.
We can introduce a piecewise function that contains separate $L_1$ and $L_2$ penalties, then set a specific threshold $\epsilon$ where one
or the other is used. To do this, we utilize a Berhu's (Reverse Huber) Loss.

Denote $y_i = |w_i - w_i^*|$. Frequently used in regressions, the Huber's Loss function takes the form
\begin{equation}
\label{huber_loss}
    h(y_i) = \begin{cases}
        \frac{y_i^2}{2\epsilon} & \text{if } y_i \leq \epsilon, \\
        y_i - \frac{\epsilon}{2} & \text{if } y_i > \epsilon.
    \end{cases}
\end{equation}
It applies a quadratic penalty for small deviations, and a linear penalty for large deviations. This mitigates the influence of
extreme outliers in regression while maintaining differentiability throughout.
Berhu's Loss, also known as Reverse Huber's Loss, inverts the order of the quadratic and linear functions, and is added to the objective function; for a turnover
$y_i$ and a threshold scalar $\epsilon > 0$, the Berhu Loss is
defined as
\begin{equation}
\label{berhu_loss}
    b(w_i) = \begin{cases}
        y_i & \text{if } y_i \leq \epsilon, \\
        \frac{y_i^2 + \epsilon^2}{2\epsilon} & \text{if } y_i > \epsilon.
    \end{cases}
\end{equation}

The Berhu function works similarly to an Elastic Net, but only applies either an $L_1$ norm or $L_2$ norm at a time. For minor weight adjustments where
$y_i \leq \epsilon$, the penalty is linear, promoting sparsity. For significant weight adjustments where $y_i \geq \epsilon$, the
penalty becomes quadratic and prevents concentrated trades.

Intuitively, Elastic Net sharpens but does not completely eliminate the smooth edges of the feasible set while the Berhu's
function shrinks the areas of pointy edges but still preserve the indifferentiable areas. Both encourage more control variables
to go towards zero upon contact with the objective function at the optimal solution, just in slightly different fashions. As a
trade increases beyond $\epsilon$, Berhu's function becomes quadratic while a comparable Elastic Net still maintains a linear
term. Therefore, Berhu's functions penalize the large outliers more aggressively than the similar Elastic Net functions.

\begin{figure}
    \centering
    \includegraphics[scale = 0.45]{huber_berhu_transparent.png}
    \caption{Berhu's Loss (Left) and Huber's Loss (Right).}
    \label{fig:loss_comparison}
\end{figure}

This function is convex and differentiable everywhere. That said, its piecewise nature makes it difficult to implement directly
into a solver, as many solvers do not accept piecewise functions natively or have a designated Berhu function.  To resolve this,
we'll create an epigraph formulation with auxiliary variables.

Define two non-negative auxiliary variables $u_i \geq 0$ and $v_i \geq 0$ subject to the following constraints:
\begin{equation*}
	u_i + v_i \geq y_i, \quad u_i \leq \epsilon,
\end{equation*}
where $y$ is defined in \eqref{eq:tvr-slack}.  The Berhu Loss then becomes
\begin{equation}
    \label{berhu_epigraph}
	b(y) = \sum_i \left(u_i + v_i + \frac{v_i^2}{2\epsilon}\right),
\end{equation}

When $y_i \leq \epsilon$, to minimize the penalty, the solver will set $v_i = 0$, since the quadratic penalty term
$\frac{v_i^2}{2\epsilon}$ grows more rapidly than the linear term $u_i$, being more "expensive" in the objective function. The
constraint $u_i + v_i \geq y_i$ forces $u_i = y_i$, making the penalty exactly $y_i$, matching the linear branch of
\eqref{berhu_loss}. When $y_i > \epsilon$, the minimizer pushes $u_i$ to its upper bound $\epsilon$, resulting in $v_i$ being
non-zero $v_i = y_i - \epsilon$, while the cost becomes $\frac{y_i^2 + \epsilon ^ 2}{2\epsilon}$. This matches the value of the
quadratic branch of \eqref{berhu_loss}. Thus, we've reformulated the piecewise Berhu function \eqref{berhu_loss} into only an
epigraph with auxiliary variables.

A penalty term based on \eqref{berhu_epigraph} can be defined as
\begin{equation}
	b_t(y) = \gamma \sum \left(u_i + v_i + \frac{v_i^2}{2\epsilon}\right) = \gamma \sum_i{\left[y_i + \frac{1}{2\epsilon}\max[y_i - \epsilon, 0]^2\right]},
\end{equation}
where $\gamma > 0$ is an adjustable exogenous scale coefficient that determines the scale of this term relative to the objective
function $q(w)$ and other penalty terms present. All terms in the objective function are either quadratic or linear, so the
function is convex.

Once again, it is important to consider the impact of changes in the parameters. Threshold $\epsilon$ is the inflection point where the
penalty transitions from linear to quadratic. A small $\epsilon$ means the optimizer switches to quadratic penalties at very small turnover values,
resulting in a more dense portfolio with fewer turnovers pushed to exactly $0$.

A large $\epsilon$ means the penalty remains linear for a wider range of turnovers, promoting sparsity in the same way the Elastic Net acted.

As stated before, the difference between \eqref{elastic_net_objective} and \eqref{berhu_epigraph} is that \eqref{berhu_epigraph} imposes only one penalty per asset,
while \eqref{elastic_net_objective} imposes both simultaneously. So in an Elastic Net, the "effective" $L_1$ penalty is always slightly increased, even when
$y_i$ is almost negligible. For the Berhu's Loss, one can set a concrete threshold $\epsilon$ past which they wish to start imposing quadratic penalties.

% ADD ADDITIONAL USAGES OF BERHU/HUBER?

\section{Huber's Loss Objective Function}
\subsection{Target Tracking}
\label{target_objective_section}

Up to this point, the base objective family $q(w)$ we used was a quadratic function along with some linear or quadratic penalty
terms represented by $p(w)$, an example of which is \eqref{gamma_matrix}. A target vector $\tilde{w}$ may represent a benchmark
portfolio or a model driven portfolio. It is often desirable to form portfolios as close as possible to such a given portfolio
subject to various constraints.

The Huber's Loss function \eqref{huber_loss} is particularly interesting as an alternative objective function for
\eqref{gamma_matrix}. For an individual tracking error $x_i = w_i - \tilde{w}_i$, and a threshold parameter $\delta$, the
Huber's Loss function can be written as
\begin{equation*}
    h(x_i) = \begin{cases}
        \frac{v_i^2}{2\delta} & \text{if } v_i \leq \delta, \\
        v_i - \frac{\delta}{2} & \text{if } v_i > \delta.
    \end{cases}
\end{equation*}
Similar to \eqref{berhu_epigraph}, this can be reformulated with auxiliary variables $u_i\in [\delta, \delta]$, $v_i^+\geq 0$,
and $v_i^-\geq 0$:
\begin{equation}
	\label{huber_loss_deviation_epigraph}
	h(x_i) = h(u_i + v_i^+ - v_i^-) = \frac{1}{2}u_i^2 + \delta (v_i^+ + v^-).
\end{equation}
The counterpart to \eqref{gamma_matrix} can be written as
\begin{equation*}
	\frac{1}{2}\sum_i{c_iu_i^2} + \delta \sum_i{c_i(v_i^+ + v_i^-)}.
\end{equation*}
Moreover, $\{\delta_i\}$ can be asset specific and the above becomes
\begin{equation*}
	\frac{1}{2}\sum_i{c_iu_i^2} + \sum_i{\delta_i c_i(v_i^+ + v_i^-)}.
\end{equation*}

CVXPY \cite{cvxpy} contains designated Huber functions that allows direct specification without slack vectors $u$, $v^+$, and
$v^-$. Behind the scene, however, it carries out the similar conversions from Huber functions.

% this part needs rewrite.
Perhaps unsurprisingly, this epigraph is the mirror opposite of our previously derived reformulation for the Berhu's Loss.
\eqref{huber_loss_deviation_epigraph} can then be made the entire objective function, along with other penalties such as alpha
gating, Elastic Net, or Berhu's Loss. In portfolio optimization, minimizing the squared tracking error (deviation from
$\tilde{w}$) or optimizing a covariance matrix $\Sigma$ can make the optimizer sensitive to outliers. If some $\tilde{w}$
predicts an extreme, inaccurate weight for a single stock, the optimizer can appropriate a comparatively massive amount of
exposure. A Huber's Loss base objective function applies an $L_2$ norm for small $|w - \tilde{w}|$, and an $L_1$ norm for larger
turnovers, preventing anomalous $\tilde{w}$ values from hijacking total exposure allocation. The value of $\lambda$ dictates how
much the solver discourages responses to large outliers in $\tilde{w}$.

% \section{Numerical Experiments}
% \subsection{contents tbd}
% \section{Conclusion}
% \subsection{Parallels with Machine Learning and Parameter Searching}
% \subsection{Next steps in research}

% --- References ---
\newpage
% Note: The reference section does not count against the 20-page limit.
\bibliographystyle{plain}
\nocite{*}
\bibliography{references}

% \begin{thebibliography}{99}

% \bibitem{author2026}
% Author, A., \& Coauthor, B. (2026). Title of foundational scientific paper. \textit{Journal of Computational Science}, 12(3), 101--115.

% \bibitem{solver2025}
% Developer, C. (2025). \textit{Convex Optimization Solvers and Algorithms}. Academic Press.

% \end{thebibliography}

\end{document}
